% ============================================================
\chapter{Descripción de la herramienta software desarrollada}
\label{cap:herramienta}
% ============================================================

El sistema kubescan se organiza en tres capas que transforman un
directorio de manifiestos \acs{YAML} en una puntuación de riesgo de cadena de
ataque, y se distribuye como una herramienta de línea de comandos instalable.
Su construcción abarca la obtención del conjunto de datos, la derivación del
grafo de clúster que representa las relaciones de seguridad entre manifiestos,
el diseño y entrenamiento de cada capa, y las prácticas de ingeniería de
software aplicadas al desarrollo. La implementación completa se encuentra
disponible en el repositorio indicado en el Anexo~\ref{ap:repositorio}.

% -----------------------------------------------------------
\section{Arquitectura general del sistema}
% -----------------------------------------------------------

El sistema kubescan implementa una arquitectura de tres capas que transforma
un directorio de manifiestos \acs{YAML} en una puntuación de riesgo de cadena de ataque
(Figura~\ref{fig:arquitectura}):

\begin{enumerate}
  \item \textbf{Capa~1 (Random Forest):} extrae 28 características por manifiesto
    (25 indicadores binarios y 3 agregados derivados) y produce un
    \textit{risk\_score} \(\in [0,1]\).
  \item \textbf{Capa~2 (\ac{GAT}):} modela el clúster como un grafo donde los nodos
    son manifiestos y las aristas representan relaciones de seguridad; produce una
    probabilidad de cadena de ataque \(\in [0,1]\).
  \item \textbf{Capa~3 (\textit{Ensemble} \acs{GA}):} combina las señales de las dos capas anteriores
    con una señal binaria de escape según la ecuación:
\end{enumerate}

\begin{equation}
  \text{score}(C) = w_{\text{rf}} \cdot \overline{\text{risk}}(C)
                  + w_{\text{gnn}} \cdot p_{\text{chain}}(C)
                  + w_{\text{esc}} \cdot \mathbf{1}[\exists\,\text{nodo de escape en }C]
  \label{eq:ensemble}
\end{equation}
\addequationtoindex{Función de puntuación del \textit{ensemble} (Capa 3)}

donde los pesos \(w_{\text{rf}}, w_{\text{gnn}}, w_{\text{esc}}\) son optimizados
mediante Algoritmo Genético sobre las predicciones \textit{out-of-fold} de la
validación cruzada (Sección~\ref{sec:capa3_diseno}).

\begin{figure}[htbp]
  \centering
  \caption{Arquitectura del sistema kubescan}
  \label{fig:arquitectura}
  \includegraphics[width=\textwidth]{fig_arquitectura}

  \fuentepropia
\end{figure}

La Figura~\ref{fig:arquitectura} muestra el proceso de tres capas, desde los
manifiestos \acs{YAML} de entrada hasta la clasificación de riesgo del clúster. El
vector de nodo del grafo de entrada de la Capa~2 reserva un atributo adicional
(índice~25) para una puntuación de severidad por manifiesto, calculada como suma
ponderada de las banderas binarias mediante \texttt{compute\_risk\_score}. En el
corpus de entrenamiento esa puntuación es una heurística determinista, no la salida
del clasificador Random Forest; la Sección~\ref{sec:construccion_grafo}
detalla su cobertura efectiva.

% -----------------------------------------------------------
\section{Construcción del conjunto de datos}
\label{sec:construccion_dataset}
% -----------------------------------------------------------

La construcción del conjunto de datos parte de la recolección de manifiestos
\acs{YAML} procedentes de fuentes públicas heterogéneas y avanza hasta la
extracción de las características que alimentan la Capa~1. El origen y la
composición de los manifiestos recopilados determinan la cobertura de patrones de
configuración errónea disponible, y el proceso de extracción de características
transforma cada manifiesto en un vector estructurado de dimensión fija. El
conjunto resultante es la base tanto del clasificador de la Capa~1 como del grafo
de clúster construido en la Sección~\ref{sec:construccion_grafo}.

\subsection{Fuentes de datos}
\label{sec:fuentes_datos}

El conjunto de datos tabular final contiene 2.827~filas y 58~columnas, combinando
cinco fuentes (agrupadas en cuatro categorías, dado que \citet{Rahman2022} aportan dos
plataformas distintas con características de distribución diferentes; ver Tabla~\ref{tab:dataset}):

\begin{table}[htbp]
\centering\small
\caption{Composición del conjunto de datos de manifiestos Kubernetes}
\label{tab:dataset}
\begin{tabular}{@{}p{5.4cm}rp{4.9cm}@{}}
\toprule
Fuente & Filas & Rol \\
\midrule
\citet{Rahman2022} — GitHub & 1.510 & Manifiestos reales, etiquetados \\
\citet{Rahman2022} — GitLab & 396 & Ídem, plataforma distinta \\
BishopFox BadPods \citep{BishopFox2021} & 128 & Pods vulnerables intencionales \\
Kubernetes Goat \citep{KubernetesGoat} & 14 & Escenarios de vulnerabilidad \\
Repositorios de seguridad adicionales & 779 & CTF, \textit{fixtures} de herramientas, labs; etiqueta derivada por regla \\
\midrule
\textbf{Total} & \textbf{2.827} & \\
\bottomrule
\end{tabular}
\fuentepropia
\end{table}

La procedencia de la etiqueta binaria no es homogénea entre las fuentes de la
Tabla~\ref{tab:dataset}, y esa heterogeneidad condiciona la interpretación de los
resultados de la Capa~1. Las 1.906~filas de \citet{Rahman2022}
conservan la anotación \textit{INSECURE} del corpus original, ajena al extractor
de características de este trabajo: solo 245 de las 1.510~filas de GitHub y 87 de
las 396 de GitLab coinciden con la disyunción de las columnas de configuración errónea
observadas. Las 779~filas de los repositorios de seguridad adicionales, en
cambio, carecen de anotación en origen y reciben una etiqueta derivada por regla
en la ingesta (\texttt{compute\_label} en
\texttt{ingest\_attack\_repos.py}), definida como la disyunción de las columnas de
configuración errónea; las 779 satisfacen por construcción esa equivalencia exacta. El
corpus final combina, por tanto, verdad de referencia externa y etiquetado por
regla, circunstancia que se retoma al analizar la circularidad entre etiqueta y
características en la Sección~\ref{sec:rf_importance}.

La incorporación de BadPods fue esencial para mitigar la escasez de ejemplos de
configuraciones erróneas críticas en los repositorios reales: antes de añadir esta fuente,
\texttt{TRUE\_HOST\_PID} aparecía en solo el 0,3\,\% de las filas; tras la incorporación,
la cobertura ascendió al 3,78\,\%. En cuanto a los conjuntos de datos de referencia
existentes para intrusión en red, como UNSW-NB15 \citep{Moustafa2015}, se estudió su
metodología de construcción como guía para el diseño del proceso de etiquetado y la
estrategia de balanceo de clases, dado que sus particiones curadas de entrenamiento y
evaluación presentan un desequilibrio de clases del mismo orden de magnitud que el de
este trabajo (56,3\,\% seguros frente a 43,7\,\% mal configurados). El conjunto GenKubeSec
\citep{Malul2024}, publicado en 2024, no estaba disponible durante el desarrollo de
este trabajo; en su lugar se empleó Checkov \citep{Checkov2021} directamente
sobre los ficheros \acs{YAML} descargados como fuente de validación multiherramienta
equivalente.

\subsection{Características extraídas}

El extractor de características aplica 28 reglas de análisis estático a cada manifiesto
\acs{YAML}, organizadas en tres grupos:

\begin{itemize}
\raggedright
  \item \textbf{18 banderas Rahman} (subconjunto del trabajo original): \texttt{TRUE\_HOST\_PID},
    \texttt{CAP\_SYS\_ADMIN}, \texttt{SEC\_CONT\_OVER\_PRIVIL}, \texttt{INSECURE\_HTTP},
    etc.
  \item \textbf{3 características derivadas}: \texttt{cap\_misuse} (OR de las dos banderas de
    capacidades de sistema), \texttt{all\_secrets} (OR de los indicadores de exposición
    de secretos), \texttt{total\_misconfigs} (suma de todas las banderas activas).
  \item \textbf{7 características extendidas}: \texttt{NO\_RUN\_AS\_NON\_ROOT},
    \texttt{IMAGE\_USES\_LATEST}, \texttt{SA\_AUTOMOUNT\_TOKEN}, \texttt{HOSTPATH\_MOUNT},
    etc. Seis de ellas presentan una procedencia mixta de valores; la séptima,
    \texttt{HOSTPATH\_MOUNT}, se extrae siempre del propio manifiesto.
\end{itemize}

La procedencia de los valores de las seis características extendidas con
equivalente en Checkov \citep{Checkov2021} —\texttt{NO\_RUN\_AS\_NON\_ROOT},
\texttt{NO\_READ\_ONLY\_ROOT\_FS}, \texttt{IMAGE\_USES\_LATEST},
\texttt{SA\_AUTOMOUNT\_TOKEN}, \texttt{USES\_DEFAULT\_SA} y
\texttt{UNTRUSTED\_REGISTRY}— se reparte en tres vías sobre las 2.827~filas del
corpus: 921~filas (32,6\,\%) reciben el valor extraído directamente por el
analizador \acs{YAML}, 1.175~filas (41,6\,\%) se rellenan con la salida de
Checkov a través de la tabla \texttt{CHECKOV\_FILL}, y las 731~restantes
(25,9\,\%) se imputan mediante la mediana de columna. Ese relleno mediante
\texttt{CHECKOV\_FILL} opera exclusivamente en la carga de la Capa~1. El constructor
del grafo lee las columnas de \texttt{FEATURE\_COLS} directamente del conjunto
tabular e imputa como 0 toda celda vacía, sin consultar las columnas
\texttt{checkov\_*}, por lo que las seis características valen 0 en el vector de
nodo de las 1.906~filas procedentes de Rahman et~al. La cobertura de
Checkov tampoco es un mecanismo general aplicable a toda fila con
\acs{YAML} disponible, sino específico de la fuente: Checkov se ejecutó únicamente sobre el subconjunto de
GitHub —1.175 de sus 1.510~filas— y no produjo salida alguna para los
repositorios de seguridad adicionales, GitLab, BadPods ni Kubernetes Goat, cuyas
filas se resuelven por el analizador \acs{YAML} o por mediana. La séptima
característica extendida queda al margen de este esquema, ya que
\texttt{HOSTPATH\_MOUNT} se deriva por completo del análisis estático del
manifiesto y está poblada para las 2.827~filas.

La Capa~2 (\ac{GNN}) consume únicamente las 18~banderas Rahman y las 7~características extendidas
(25 en total, \texttt{FEATURE\_COLS} en \texttt{yaml\_parser.py}) como vector de nodo,
excluyendo las 3~características derivadas —agregados calculados a partir de las otras 25 que
no aportan información estructural adicional al grafo—; ver
Sección~\ref{sec:features_table}. De estas 25, tres presentan varianza casi nula pero
se mantienen en el entrenamiento, con importancia residual o nula
(Sección~\ref{sec:rf_importance}): \texttt{SECCOMP\_UNCONFINED} (15/2.827 activa),
\texttt{VALID\_TAINT\_SECRET} (0/2.827 activa) y \texttt{NO\_NETWORK\_POLICY}
(2.818/2.827 activa, sin poder discriminativo). Las importancias de característica
registradas por el modelo Random Forest confirman ese aporte marginal:
\(1{,}6\times10^{-5}\) para \texttt{SECCOMP\_UNCONFINED}, \(0{,}0\) para
\texttt{VALID\_TAINT\_SECRET} y \(2{,}4\times10^{-5}\) para
\texttt{NO\_NETWORK\_POLICY}. Se conservan para preservar la estabilidad del
contrato de características (\texttt{FEATURE\_COLS}) y la correspondencia exacta
con el conjunto de banderas de \citet{Rahman2022}; no se ha realizado
un estudio de ablación que cuantifique el efecto de su eliminación.

% -----------------------------------------------------------
\section{Construcción del grafo de clúster}
\label{sec:construccion_grafo}
% -----------------------------------------------------------

Cada directorio de manifiestos se representa, mediante NetworkX
\citep{Hagberg2008}, como un grafo dirigido
\(G=(V, E)\) donde \(|V|\) es el número de manifiestos del clúster y
\(E\) contiene aristas de cinco tipos semánticamente distintos
(Tabla~\ref{tab:edge_types}). La Figura~\ref{fig:grafo_cluster} ilustra un ejemplo
con cuatro nodos que incluye los tres tipos de nodo posibles: escape, movimiento
lateral y limpio. Un \textit{namespace} o espacio de nombres es la partición
lógica con la que Kubernetes agrupa y aísla recursos dentro de un mismo clúster,
y delimita el alcance de tres de los cinco tipos de arista. Las reglas de
asignación de cada tipo de nodo y de arista se detallan en las dos subsecciones
siguientes.

\begin{table}[htbp]
\centering\small
\caption{Tipos de aristas en el grafo de clúster}
\label{tab:edge_types}
\begin{tabular}{@{}p{3cm}cp{3.6cm}p{3.8cm}@{}}
\toprule
Tipo & Código & Dirección & Criterio \\
\midrule
Proximidad de directorio & 0 & Bidireccional & Mismos dos componentes finales de ruta \acs{YAML} \\
Alcance de privilegio & 1 & Dirigida (escape\(\to\)mismo \textit{namespace}) & Bandera de escape activa \\
Movimiento lateral \acs{SA} & 2 & Dirigida (lateral\(\to\)mismo \textit{namespace}) & Bandera de movimiento lateral \\
Co-espacio de nombres (\texttt{SEMANTIC\_NS}) & 3 & Bidireccional & Mismo \textit{namespace} K8s \\
Escalada \ac{RBAC} & 4 & Dirigida (elevado\(\to\)mismo \textit{namespace}) & \ac{SA} vinculado a rol privilegiado \\
\bottomrule
\end{tabular}
\fuentepropia
\end{table}

Cada nodo \(v_i\) lleva un vector de características \(\mathbf{x}_i \in \mathbb{R}^{26}\):
los índices 0–24 corresponden a las 25 características binarias de la Capa~1 y el
índice~25 contiene una puntuación de severidad por manifiesto. Esa puntuación
procede de \texttt{compute\_risk\_score}, una suma ponderada de las mismas banderas
binarias normalizada a [0,1], y no de una predicción del clasificador
Random Forest. Está poblada únicamente para los 82~clústeres procedentes
del corpus de Rahman et~al.; en los 517~restantes la columna llega vacía al
constructor del grafo y se imputa como 0, de modo que el índice~25 no aporta señal
en el 86\,\% de los grafos.

\begin{figure}[htbp]
  \centering
  \caption{Ejemplo de grafo de clúster Kubernetes con cuatro nodos}
  \label{fig:grafo_cluster}
  \includegraphics[width=0.72\textwidth]{fig_grafo_cluster}

  \fuentepropia
\end{figure}

La Figura~\ref{fig:grafo_cluster} muestra los tres tipos de nodo (escape, movimiento
lateral y limpio) y las aristas dirigidas que cada uno origina. El nodo rojo (pod-A) tiene
banderas de escape activas y emite aristas de tipo~1 (\textit{privilege\_reach}) hacia
todos los demás nodos de su mismo \textit{namespace}. El nodo naranja (pod-B) tiene banderas de
movimiento lateral y emite aristas de tipo~2 (\textit{sa\_lateral}). La arista de
tipo~4 (\ac{RBAC}) representa la vinculación de un ServiceAccount a un rol
privilegiado.

El conjunto de datos de grafos final comprende 599~grafos originales (25~clean, 537~isolated,
37~attack\_chain) y 1.154~grafos totales tras la aumentación de datos, cuyas variantes
se destinan exclusivamente al conjunto de entrenamiento. De los 555~grafos aumentados
que componen ese \textit{pool}, solo 390 entran efectivamente en el entrenamiento: los
165~restantes derivan de clústeres asignados a validación o test y quedan excluidos por
la regla de asignación consciente de grupos descrita en la
Sección~\ref{sec:aumentacion}. Esta cifra incorpora el corpus
completo de \citet{Rahman2022}, cuyos \textit{namespaces} cruzados
se resolvieron mediante descarga autenticada desde GitHub, lo que amplió
sustancialmente el número de clústeres \textit{isolated} respecto a las fases
experimentales preliminares descritas en la Sección~\ref{sec:layer2} del
Capítulo~\ref{cap:evaluacion}.

\subsection{Detección de escape y movimiento lateral}

{\raggedright
La clasificación de nodos utiliza dos conjuntos de banderas definidos en el código
como constantes inmutables (\texttt{frozenset}). Un nodo se clasifica como de
\textbf{escape} si tiene activa al menos una de las siguientes 8~banderas:
\texttt{TRUE\_HOST\_PID}, \texttt{TRUE\_HOST\_IPC}, \texttt{TRUE\_HOST\_NET},
\texttt{DOCKERSOCK\_PATH}, \texttt{CAP\_SYS\_ADMIN}, \texttt{CAP\_SYS\_MODULE},
\texttt{SEC\_CONT\_OVER\_PRIVIL} o \texttt{HOSTPATH\_MOUNT}. Este último vector fue
incorporado durante el desarrollo al detectar que \textit{badpods\_hostpath} —que monta
el sistema de ficheros del anfitrión— era incorrectamente etiquetado como
\textit{isolated} al no tener ninguna otra bandera de escape activa. Con la corrección,
el número de clústeres de cadena de ataque del corpus preliminar ascendió de 13
a~15; la trayectoria completa del tamaño del corpus y del número de cadenas a lo
largo de las fases experimentales se recoge en la
Tabla~\ref{tab:gnn_evolution}.\par}

Un nodo se clasifica como de \textbf{movimiento lateral} si tiene activa al menos una
de las siguientes 4~banderas: \texttt{SA\_AUTOMOUNT\_TOKEN} (\textit{token} de cuenta de servicio
automontado), \texttt{USES\_DEFAULT\_SA} (usa la cuenta de servicio por defecto del
\textit{namespace}), \texttt{WITHIN\_MANIFEST\_SECRET} (secreto referenciado en el
manifiesto) o \texttt{ALLOW\_PRIVI} (contenedor con capacidad de escalada). Las aristas
de movimiento lateral (tipo~2) se construyen de forma análoga a las de escape: se emite
una arista dirigida desde cada nodo lateral hacia los demás nodos de su mismo
\textit{namespace}.

La convivencia de los cinco tipos sobre un mismo par de nodos se resuelve por orden de
inserción y no por una jerarquía explícita de prioridades. Las aristas se añaden en la
secuencia proximidad de directorio (tipo~0), alcance de privilegio (tipo~1), movimiento
lateral (tipo~2), coespacio de nombres (tipo~3) y escalada \ac{RBAC} (tipo~4). Las de
alcance de privilegio son las únicas que sobrescriben incondicionalmente cualquier
etiqueta previa, de modo que un nodo con capacidad de escape nunca queda degradado a
mero vecino; los cuatro tipos restantes se emiten bajo la guarda \texttt{not
G.has\_edge}, por lo que una arista de proximidad de directorio ya presente impide que
el par se reetiquete como movimiento lateral, y las aristas \ac{RBAC}, al añadirse en
último lugar, nunca desplazan a las de movimiento lateral.

\subsection{Detección de privilegio \acs{RBAC}}

La construcción de aristas \ac{RBAC} (tipo~4) requiere dos pasadas sobre el conjunto de
manifiestos. La primera pasada identifica los roles privilegiados: aquellos cuyo nombre
está en el conjunto \{cluster-admin, admin, edit\} o cuyas reglas contienen verbos de
escalada \{*, escalate, bind, impersonate\}. La segunda pasada recorre los objetos
RoleBinding y ClusterRoleBinding: solo cuando el campo
\textit{roleRef} referencia un rol privilegiado identificado en la primera pasada se
emite una arista \ac{RBAC} desde el pod que usa esa cuenta de servicio hacia los demás
nodos de su mismo \textit{namespace}, coherentemente con el ámbito de los objetos
RoleBinding de los que se derivan las cuentas elevadas.

% -----------------------------------------------------------
\section{Capa 1: Random Forest}
\label{sec:capa1_diseno}
% -----------------------------------------------------------

La Capa~1 clasifica cada manifiesto de forma individual a partir de las
características descritas en la Sección~\ref{sec:construccion_dataset}. Se entrenan dos modelos
independientes sobre el mismo conjunto de características: un clasificador
binario que distingue manifiestos seguros de mal configurados, y un modelo de
severidad que estima el grado de gravedad de la configuración incorrecta
detectada. Ambos comparten el conjunto de 28~características y los mismos
hiperparámetros de bosque, y difieren únicamente en la variable objetivo.

\subsection{Modelo binario}

Se entrena un \texttt{RandomForestClassifier} \citep{Breiman2001} de scikit-learn
\citep{Pedregosa2011} con los hiperparámetros siguientes:

\begin{itemize}
  \item \texttt{n\_estimators = 500}
  \item \texttt{max\_features = sqrt} (clasificación estándar)
  \item \texttt{min\_samples\_leaf = 2} (evita hojas de muestra única)
  \item \texttt{class\_weight = balanced} (compensa ratio 1,29:1)
  \item \texttt{random\_state = 42}
\end{itemize}

El modelo se entrena sobre 2.261~muestras (80\,\% del total) y se evalúa en 566
muestras de test no vistas durante el entrenamiento. La probabilidad predicha para la
clase misconfigured (clase~1) se usa como \textit{risk\_score} del nodo en el grafo.

\subsection{Modelo de severidad}

Además del modelo binario, se entrena un clasificador de severidad de tres clases
(\textit{clean}, \textit{low-medium}, \textit{high-critical}) para ofrecer información
más granular al operador de seguridad. Este modelo no interviene en el cálculo del
grafo; se incluye como salida adicional de la Capa~1.

% -----------------------------------------------------------
\section{Capa 2: Graph Attention Network}
\label{sec:capa2_diseno}
% -----------------------------------------------------------

La Capa~2 opera sobre el grafo de clúster construido en la
Sección~\ref{sec:construccion_grafo} y clasifica el clúster completo según el
riesgo de cadena de ataque que presenta. Su diseño se apoya en tres decisiones
acopladas: una arquitectura \acs{GAT} con \textit{embedding} de tipo de arista, un
procedimiento de entrenamiento con validación cruzada estratificada de
5~pliegues, y una estrategia de aumentación de datos que mitiga la escasez de
clústeres con cadena de ataque confirmada.

\subsection{Arquitectura}

El modelo \ac{GAT} \citep{Velickovic2018} implementado tiene la siguiente estructura,
con las dimensiones exactas verificadas en \texttt{gat\_encoder.py}:

\begin{enumerate}
\raggedright
  \item \textbf{\textit{Embedding} de tipo de arista}: \texttt{Embedding(5, 8)}.
        Los cinco tipos de arista (proximidad de directorio, alcance de
        privilegio, lateral por ServiceAccount, espacio de nombres y
        \ac{RBAC}) son categóricos: codificarlos como un escalar continuo
        implicaría un orden espurio (el tipo~4 «valdría» cuatro veces el
        tipo~1). Cada tipo se proyecta a un vector aprendido de 8~dims que
        \texttt{GATConv} consume como atributo de arista
        (\textit{edge\_dim=8}).
  \item \textbf{Proyección inicial}: \texttt{Linear(26, 256)} \(+\) \texttt{LayerNorm(256)}
        \(+\) \ac{ReLU} \(+\) \texttt{Dropout(0.3)}.
        Proyecta el vector de nodo de 26~dims a 256~dims
        (\(64 \times 4\,\text{cabezas}\)).
  \item \textbf{Capas \ac{GAT}~0 y~1} (\textit{concat=True}):
        \texttt{GATConv(256, 64, heads=4)} con \textit{edge\_dim=8};
        salida 256~dims \((64 \times 4)\); seguidas de \texttt{LayerNorm}, \ac{ELU},
        conexión residual y \texttt{Dropout(0.3)}.
  \item \textbf{Capa \ac{GAT}~2} (\textit{concat=False, heads=1}):
        \texttt{GATConv(256, 64, heads=1)}; salida 64~dims
        (una sola cabeza de atención que produce directamente
        el \textit{embedding} final del nodo, sin concatenación);
        seguida de \texttt{LayerNorm}, \acs{ELU} y \texttt{Dropout(0.3)}
        (sin conexión residual: las dimensiones de entrada y salida difieren,
        256 frente a 64).
  \item \textbf{Global \textit{pooling}}: \texttt{global\_mean\_pool} y \texttt{global\_max\_pool}
        sobre los 64-dim \textit{embeddings} de todos los nodos; concatenación
        \(\to\) 128~dims (\(64 + 64\)).
  \item \textbf{Clasificador \ac{MLP}}: \texttt{Linear(128, 64)} \(+\) \acs{ReLU} \(+\)
        \texttt{Dropout(0.3)} \(+\) \texttt{Linear(64, 3)}.
\end{enumerate}

La elección de \ac{GAT} sobre \ac{GCN} responde a que el mecanismo de atención permite al modelo
aprender qué aristas son más informativas para la señal de ataque: una arista de
\textit{privilege\_reach} debe recibir mayor peso que una arista de proximidad de
directorio. El \textit{pooling} combinado de media y máximo captura dos señales
complementarias: la media refleja la postura de seguridad promedio del clúster,
mientras que el máximo identifica el nodo más peligroso, ambas necesarias para detectar
cadenas de ataque que pueden estar concentradas en un único nodo crítico.

\subsection{Entrenamiento}

El modelo se entrena con PyTorch Geometric \citep{Fey2019}, construido
sobre PyTorch \citep{Paszke2019}, bajo los siguientes hiperparámetros: optimizador AdamW con tasa de aprendizaje
\(5\times10^{-4}\) y \textit{weight\_decay}\,=\,\(10^{-4}\);
\textit{scheduler} CosineAnnealingLR con T\_max\,=\,300 y
\(\eta_{\min}\,=\,10^{-5}\); función de pérdida CrossEntropyLoss con pesos
de clase calculados por frecuencia inversa;
\textit{gradient clipping} con norma máxima 2,0; \textit{early stopping} con
paciencia de 50 épocas; tamaño de lote 8; semilla global 42 (idéntica para el
entrenamiento y para la generación de particiones); y validación cruzada
estratificada de 5~pliegues. El \textit{checkpoint} de cada pliegue se selecciona
por la mejor Precisión@5 de validación, con la F1 macro como criterio de
desempate (\texttt{--select-by p5}), alineando la selección de modelos con la
métrica primaria del sistema en lugar de con la métrica de clasificación.

El guion de entrenamiento (\texttt{train\_gnn.py}) expone tres opciones que
seleccionan variantes de esa configuración y que hacen reproducibles los estudios
de ablación recogidos en el Capítulo~\ref{cap:evaluacion}. La opción
\texttt{-{}-conv} escoge el operador de convolución sobre el grafo entre
\texttt{gat}, la \ac{GAT} con \textit{embedding} de tipo de arista descrita en la
arquitectura y empleado en la configuración desplegada, y \texttt{gcn}, una
\ac{GCN} de referencia que ignora los tipos de arista y pondera todos los vecinos
por igual. La opción \texttt{-{}-focal-loss}, acompañada del parámetro de enfoque
\texttt{-{}-focal-gamma} (valor 2 por defecto), sustituye la entropía cruzada
ponderada por la \textit{focal loss} \citep{Lin2017}, que reduce el peso de los
ejemplos ya bien clasificados para concentrar el gradiente en los difíciles. La
opción \texttt{-{}-select-by} fija el criterio de selección del \textit{checkpoint}
de cada pliegue entre \texttt{f1} (F1 macro de validación) y \texttt{p5}
(Precisión@5 de validación), con la métrica restante como criterio de desempate.
La configuración desplegada corresponde a \texttt{-{}-conv gat}, sin
\texttt{-{}-focal-loss} y con \texttt{-{}-select-by p5}; las alternativas se
cuantifican en las Secciones~\ref{sec:ablacion_perdida}
y~\ref{sec:ablacion_arquitectura}.

Los pesos de clase se calculan en cada pliegue de entrenamiento mediante la
siguiente fórmula implementada en \texttt{train\_gnn.py}:

\begin{equation}
  w_i = \frac{1/n_i}{\displaystyle\sum_{j=0}^{K-1} 1/n_j} \cdot K
  \label{eq:class_weights}
\end{equation}
\addequationtoindex{Cálculo de pesos de clase para el entrenamiento de la GAT}

donde \(n_i\) es el número de grafos de entrenamiento de la clase \(i\) y \(K=3\)
es el número de clases. Esta fórmula normaliza los pesos para que su media sea~1,0,
proporcionalmente equivalente al comportamiento de \texttt{class\_weight=balanced} de
scikit-learn (misma ponderación relativa entre clases). Cabe destacar que los pesos se calculan
\textbf{después} de la aumentación de datos: como la aumentación genera 15~variantes
por grafo de cadena de ataque, la clase \textit{attack\_chain} pasa de ser la
minoría a ser la mayoría en el conjunto de entrenamiento. Los tres recuentos varían de
un pliegue a otro, ya que las variantes de los clústeres reservados para validación se
excluyen en cada iteración: entre 360 y 405~grafos de \textit{attack\_chain} aumentados
frente a 17--18 \textit{clean} y 361--371 \textit{isolated}. La
Ecuación~\ref{eq:class_weights} se autoadapta a esta nueva distribución
asignando mayor peso a \textit{clean} e \textit{isolated}, y menor peso a
\textit{attack\_chain}, compensando el efecto de la sobrerrepresentación creada por
la aumentación.

\subsection{Aumentación de datos}
\label{sec:aumentacion}

Para mitigar la escasez de ejemplos de cadena de ataque (37 de 599 grafos
originales), se aplican tres estrategias de aumentación exclusivamente sobre los
grafos de entrenamiento:

\begin{itemize}
  \item \textbf{\textit{Edge dropout}.} Eliminación aleatoria de entre el 10\,\%
    y el 35\,\% de aristas de los tipos~0 y~3, los menos críticos para la señal de
    ataque, con tasa muestreada uniformemente por variante.

  \item \textbf{\textit{Feature masking}.} Puesta a cero con probabilidad por
    nodo entre el 10\,\% y el 25\,\%, restringida a seis características ajenas a
    la regla de etiquetado de cadena: nunca alcanza las banderas de escape,
    \texttt{HOSTPATH\_MOUNT} ni el \textit{risk\_score}.

  \item \textbf{\textit{Subgraph sampling}.} Extracción de un subgrafo inducido
    por entre el 40\,\% y el 70\,\% de nodos, aplicada a las últimas cinco de las
    15~variantes y solo a grafos de más de 50~nodos, con \textit{BFS} sembrada en
    el nodo de mayor \textit{risk\_score}. Cada subgrafo se valida contra la regla
    de etiquetado de cadena y, si deja de cumplirla, se descarta en favor de una
    variante sobre el grafo completo, lo que evita introducir ruido de etiqueta.
\end{itemize}

Cada grafo de cadena de ataque origina 15~variantes aumentadas, lo que produce
555~grafos adicionales (37~originales \(\times\) 15~variantes) en el
\textit{pool} de aumentación. De estos, únicamente 390 —los correspondientes a
los 26~grafos de cadena cuya base pertenece a la partición de entrenamiento—
entran efectivamente en el entrenamiento, para un total de 1.154~grafos con
aumentación.

La asignación de variantes a particiones es \textbf{consciente de grupos}
(\textit{group-aware}): una variante aumentada solo puede entrar en una partición
de entrenamiento si su grafo base pertenece a esa misma partición. Las variantes
derivadas de clústeres de validación o de test se excluyen del entrenamiento por
completo, ya que entrenar con casi duplicados de un grafo de evaluación
constituiría fuga de datos.

El particionado es además \textbf{consciente de familias de plantilla}: los
clústeres originales que comparten una misma plantilla base (por ejemplo, las
ocho variantes \texttt{badpods\_*}, que difieren únicamente en qué bandera de
escape activan —salvo \texttt{badpods\_nothing-allowed}, que no activa ninguna—,
o las nueve \texttt{datadog\_*}) son casi duplicados entre sí, no
muestras independientes. Cuando una familia contiene etiquetas mixtas, sus
miembros se asignan como una unidad atómica a una única partición —nunca se
reparten entre entrenamiento, validación y test ni entre pliegues— y la
asignación emplea un reparto voraz ponderado por el número de grafos de cada
etiqueta, de modo que la proporción real de cada clase por partición siga las
fracciones solicitadas. Sin esta regla, un modelo entrenado con cinco hermanos de
una familia resuelve trivialmente al sexto en test, inflando las métricas. La
partición resultante (semilla 42) es de 815~grafos de entrenamiento (con
aumentación), 88 de validación y 86 de test. Adicionalmente, los 86~clústeres de
test se mantienen fuera de los pliegues de validación cruzada: los pliegues
particionan únicamente los 513~grafos originales restantes (32 de cadena), de
modo que ni los modelos de pliegue ni las predicciones \textit{out-of-fold}
empleadas para ajustar los pesos del \textit{ensemble} (Capa~3) observan ningún
clúster de test, directa ni indirectamente.

% -----------------------------------------------------------
\section{Capa 3: \textit{Ensemble} con Algoritmo Genético}
\label{sec:capa3_diseno}
% -----------------------------------------------------------

Los pesos \((w_{\text{rf}}, w_{\text{gnn}}, w_{\text{esc}})\) de la
Ecuación~\ref{eq:ensemble} son optimizados mediante un Algoritmo Genético
\citep{Holland1975} con 60~individuos codificados como pares
\((w_{\text{gnn}}, w_{\text{esc}})\) \(\in [0,1]^2\), con \(w_{\text{rf}}\) derivado
como \(1 - w_{\text{gnn}} - w_{\text{esc}}\) (la representación 2D garantiza que la
suma de los tres pesos sea siempre~1). La función objetivo combina la métrica de
\textit{ranking} Precisión@5 y la tasa de falsos positivos sobre clústeres
limpios \ac{FPR}\_clean, ambas definidas en la Sección~\ref{sec:metricas}, junto
con un término de regularización:

\begin{equation}
  F = 0{,}7 \cdot P@5 + 0{,}3 \cdot (1 - \text{FPR\_clean})
      - \lambda \sum_{i} \max(0,\ w_{\min} - w_i)
  \label{eq:ga_fitness}
\end{equation}
\addequationtoindex{Función objetivo (fitness) del Algoritmo Genético}

El último término es una penalización de suelo con \(w_{\min}\,=\,0{,}1\) y
\(\lambda\,=\,2{,}0\): sin ella, el optimizador es libre de anular cualquiera de
los tres pesos cuando eso puntúa marginalmente mejor sobre la muestra
\textit{out-of-fold}, lo que produce soluciones degeneradas que no generalizan al
conjunto de test. La penalización hace que ignorar una señal por completo esté
estrictamente dominado salvo que la ganancia observada supere \(\lambda\) veces el
déficit.

La selección se realiza mediante torneo de tamaño~3, con elitismo del 10\,\%; el cruce
es aritmético (media de los dos padres) con probabilidad 0,7; la mutación combina
perturbación gaussiana (\(\sigma\,=\,0{,}08\), tasa del 15\,\%) con una mutación de
frontera (tasa del 5\,\%) que salta a los vértices del símplex, seguida de
reproyección al dominio válido (\(w_{\text{rf}} \geq 0\)); la población inicial se
siembra con los vértices y puntos medios del símplex y con muestreo de Dirichlet, y el
criterio de parada es 150~generaciones.
Como \textit{baseline} determinista, antes de la optimización genética se
ejecuta una búsqueda en rejilla 2D con paso \(1/20\) (231~evaluaciones); la
comparación de su óptimo con el resultado del \ac{GA} se discute en la
Sección~\ref{sec:layer3}.

La señal de escape se define como un indicador binario que vale~1,0 si al menos un
nodo del clúster tiene alguna bandera de escape activa, y~0,0 en caso contrario. Esta
formulación evita la dilución que produciría una fracción cuando la mayoría de los
nodos del clúster son limpios.

La asignación de clase final a partir de la puntuación \textit{ensemble\_score}
se realiza mediante los siguientes umbrales fijos, definidos en
\texttt{ga\_ensemble.py}:
\begin{itemize}
  \item \textit{ensemble\_score}\,\(\geq\)\,0,60 \(\ \Rightarrow\) \texttt{ATTACK\_CHAIN}
  \item \textit{ensemble\_score}\,\(\geq\)\,0,30 \(\ \Rightarrow\) \texttt{ISOLATED\_MISCONFIG}
  \item De lo contrario \(\Rightarrow\) \texttt{CLEAN}
\end{itemize}
Estos umbrales son fijos y no se reoptimizan durante el entrenamiento del \ac{GA},
que optimiza únicamente los pesos \(w_i\) de la Ecuación~\ref{eq:ensemble}.
Permiten al operador interpretar directamente la puntuación: cualquier clúster con
\textit{score}\,\(\geq\)\,0,60 requiere revisión inmediata.

\subsection{Restricción de viabilidad estructural}
\label{sec:gate_estructural}

Una cadena de ataque multisalto requiere, por definición, alcanzabilidad entre al
menos dos manifiestos; un grafo de clúster con un único nodo no puede albergarla.
Esta propiedad, verificada sobre el corpus completo (los 592~grafos de cadena
—37 originales y 555 variantes aumentadas— tienen \(\geq\)2 nodos y \(\geq\)1 arista,
mientras que 499 de los 537 grafos
\textit{isolated} son de un solo nodo), se incorpora como restricción de
inferencia en \texttt{ga\_ensemble.py} (\texttt{MIN\_CHAIN\_NODES},
\texttt{is\_chain\_feasible}, \texttt{chain\_rank\_key}): en el
\textit{ranking} por riesgo de cadena, ningún clúster de un solo manifiesto puede
superar a un clúster estructuralmente viable, y el veredicto de un clúster
inviable se limita a \texttt{ISOLATED\_MISCONFIG} aunque su puntuación supere el
umbral de 0,60. La restricción se aplica únicamente en inferencia y evaluación;
el ajuste de los pesos del \ac{GA} se realiza sobre el \textit{ranking} sin
restricción, donde los pesos están mejor determinados (la Sección~\ref{sec:layer3}
cuantifica el efecto de esta regla sobre las métricas de test).

% -----------------------------------------------------------
\section{Implementación de kubescan}
\label{sec:implementacion}
% -----------------------------------------------------------

El sistema se distribuye como un paquete Python instalable mediante
\texttt{pip install -e kubescan/}. La estructura del paquete sigue la convención
\textit{src-layout} de \texttt{setuptools}: el código fuente reside en
\texttt{kubescan/src/kubescan/} y los módulos se organizan en tres submódulos:
\texttt{model/} (clasificadores Random Forest, \ac{GAT} y ensamblador \ac{GA}), \texttt{utils/}
(\textit{parser} \acs{YAML} y constructor de grafos) y \texttt{cli.py} (interfaz de
línea de comandos basada en Click).

El punto de entrada principal expone dos comandos. El primero,
\texttt{kubescan scan <dir>}, analiza un directorio local de manifiestos \acs{YAML} y
emite un informe de riesgo en formato texto o \ac{JSON}. El segundo,
\texttt{kubescan live}, consulta la \ac{API} de Kubernetes mediante
\texttt{kubectl}, vuelca los recursos de \textit{workloads} y \ac{RBAC} en un directorio
temporal y ejecuta el \textit{pipeline} completo sobre ellos. La resolución de los
pesos del modelo sigue la cadena: argumento \texttt{--checkpoints-dir}
\(\to\) variable de entorno \texttt{KUBESCAN\_CHECKPOINTS} \(\to\) enlace simbólico
\texttt{kubescan/checkpoints/trained/} \(\to\) \texttt{CheckpointNotFoundError},
excepción tipada derivada de \texttt{KubescanError} y definida en
\texttt{exceptions.py}.

Ambos comandos comparten el conjunto de opciones recogido en la
Tabla~\ref{tab:cli_options} del Anexo~\ref{ap:guia_uso}, definido con
Click en \texttt{cli.py}. La
opción \texttt{-{}-format} selecciona entre el informe en texto y la salida
\ac{JSON}; \texttt{-{}-show-nodes} añade el desglose por manifiesto al informe en
texto; \texttt{-{}-cluster-name} fija el nombre del clúster que encabeza el
informe, que por defecto es el del directorio analizado; y
\texttt{-{}-checkpoints-dir} altera el primer eslabón de la cadena de resolución de
pesos. Las opciones \texttt{-{}-namespace} y \texttt{-{}-all-namespaces} son
exclusivas de \texttt{kubescan live} y delimitan el ámbito de la consulta al
\ac{API} de Kubernetes. El grupo raíz añade, por su parte,
\texttt{-{}-verbose}, que eleva el nivel del registro de \texttt{WARNING} a
\texttt{DEBUG}, y \texttt{-{}-version}, que emite la versión instalada del paquete.

El informe en texto se estructura en tres bloques. El encabezado identifica el
clúster y la ruta analizada. El cuerpo emite el veredicto
(\texttt{CLEAN}, \texttt{ISOLATED\_MISCONFIG} o \texttt{ATTACK\_CHAIN}), seguido de
las señales que lo sustentan —\textit{ensemble score}, probabilidad de cadena del
\textit{ensemble} de cinco pliegues, probabilidad de clúster limpio, riesgo medio del Random Forest,
fracción de manifiestos con capacidad de escape y con capacidad de movimiento
lateral— y de los tres pesos \(w_{\text{rf}}\), \(w_{\text{gnn}}\) y
\(w_{\text{esc}}\) efectivamente aplicados. El bloque final, condicionado a
\texttt{-{}-show-nodes}, tabula los manifiestos ordenados por
\textit{risk\_score} descendente con su marca de escape (\texttt{ESC}) o
movimiento lateral (\texttt{LAT}) y sus primeras cuatro banderas activas. La
Figura~\ref{fig:cli_pipeline} del Anexo~\ref{ap:guia_uso} traza el recorrido
completo de una invocación de \texttt{kubescan scan} y reproduce una
transcripción íntegra de este informe.

La salida \ac{JSON} expone la misma información en un objeto de un solo nivel
apto para consumo automatizado: los campos \texttt{cluster},
\texttt{cluster\_dir} y \texttt{verdict} identifican el análisis; los campos
\texttt{ensemble\_score}, \texttt{chain\_probability},
\texttt{clean\_probability}, \texttt{mean\_rf\_risk} y \texttt{escape\_fraction}
recogen las señales numéricas redondeadas a seis decimales; los recuentos
\texttt{n\_manifests}, \texttt{n\_escape\_capable} y \texttt{n\_lateral\_capable}
resumen la composición del clúster; el objeto anidado \texttt{weights} contiene
los tres pesos del \textit{ensemble}; y la lista \texttt{manifests}, ordenada por
\textit{risk\_score} descendente, aporta por cada manifiesto su nombre de fichero,
su \textit{risk\_score}, los indicadores booleanos \texttt{escape\_capable} y
\texttt{lateral\_capable}, y la lista completa \texttt{flags} de características
de seguridad activas. Estos dos últimos campos son la evidencia de trazabilidad
que verifica el requisito RNF-1 en la
Sección~\ref{sec:verificacion_requisitos}.

% -----------------------------------------------------------
\section{Tabla de características del sistema}
\label{sec:features_table}
% -----------------------------------------------------------

El catálogo completo de las 28~características, con su grupo de procedencia y su
papel en el grafo —detección de escape (ESC) o de movimiento lateral (LAT)—, se
recoge en las Tablas~\ref{tab:features_a}, \ref{tab:features_a2}
y~\ref{tab:features_b} del Anexo~\ref{ap:repositorio}. De las 28, las 25 que no
son derivadas (18~Rahman + 7~extendidas) forman también el vector de nodo de la
Capa~2.

Las 8~banderas de escape (ESC, \texttt{ESCAPE\_FLAGS} en \texttt{graph\_builder.py}) son
\identcode{TRUE\_HOST\_PID}, \identcode{TRUE\_HOST\_IPC}, \identcode{TRUE\_HOST\_NET},
\identcode{DOCKERSOCK\_PATH}, \identcode{CAP\_SYS\_ADMIN}, \identcode{CAP\_SYS\_MODULE},
\identcode{SEC\_CONT\_OVER\_PRIVIL} y \identcode{HOSTPATH\_MOUNT}. Las 4~banderas de movimiento
lateral (LAT, \texttt{LATERAL\_FLAGS}) son \identcode{SA\_AUTOMOUNT\_TOKEN},
\identcode{USES\_DEFAULT\_SA}, \identcode{WITHIN\_MANIFEST\_SECRET} y \identcode{ALLOW\_PRIVI}.
La procedencia de las 7~características extendidas —mixta en seis de ellas y
siempre del propio manifiesto en \identcode{HOSTPATH\_MOUNT}— se detalla en la
Sección~\ref{sec:construccion_dataset}.

% -----------------------------------------------------------
\section{Calidad de software y proceso de desarrollo}
\label{sec:calidad}
% -----------------------------------------------------------

Cuatro herramientas integradas en un
\textit{pipeline} de pre-commit y \acs{CICD} comprueban de forma automática el
estilo, los tipos, el comportamiento y la complejidad del código, y bloquean la
integración cuando alguna comprobación falla: (1) \texttt{ruff}, linter y
formateador con reglas de estilo y seguridad; (2) \texttt{mypy} en modo estricto
(\texttt{mypy kubescan/src -{}-strict}), para el tipado estático de todo el
paquete; (3) \texttt{pytest}, para las pruebas unitarias y de integración; y
(4) \texttt{radon}, para la complejidad ciclomática. El reparto entre ambas
etapas no es uniforme: los \textit{hooks} de pre-commit ejecutan únicamente
\texttt{ruff}, \texttt{ruff-format} y comprobaciones de higiene del repositorio
(espacios finales, fin de fichero, validez de los \acs{YAML} de configuración y
tamaño máximo de los ficheros añadidos), mientras que \texttt{mypy},
\texttt{pytest} y \texttt{radon} se ejecutan en el flujo de integración continua
de GitHub Actions sobre cada \textit{push} y cada \textit{pull request}.

{\raggedright
La suite comprende siete módulos unitarios —analizador \acs{YAML}, constructor
del grafo, codificador \ac{GAT}, clasificador Random Forest, ensamblador \ac{GA},
utilidades de dispositivo y un módulo de pruebas basadas en propiedades con
Hypothesis— y un módulo de integración de extremo a extremo, cuya
verificación funcional se aporta como evidencia de RF-1 y RF-2 en la
Sección~\ref{sec:verificacion_requisitos}. La integración continua ejecuta la
suite con medición de cobertura (\texttt{-{}-cov}). Tres decisiones de diseño
sostienen la reproducibilidad y la seguridad de la carga: \texttt{set\_global\_seed}
centraliza las semillas de \texttt{random}, NumPy y PyTorch;
\texttt{resolve\_device} fija de forma determinista el \textit{backend} de
cómputo; y la carga del clasificador valida los tipos declarados por el fichero
\texttt{skops} contra una lista blanca restringida a scikit-learn y
NumPy, de modo que un \textit{checkpoint} manipulado no pueda introducir
tipos arbitrarios. En el lado de investigación, cada \acs{JSON} de resultados
incorpora un bloque de procedencia con el \textit{commit}, las semillas y las
versiones de los paquetes que lo produjeron.\par}
